% Bagian: first-order-logic
% Bab: axiomatic-deduction
% Subbagian: provability-quantifiers

% Verifikasi sifat-sifat keterderivasian yang diperlukan untuk himpunan
% konsisten maksimal dalam bab kelengkapan.

\documentclass[../../../include/open-logic-section]{subfiles}

\begin{document}

\olfileid[id]{fol}{axd}{qpr}

\olsection{\usetoken{S}{derivability} dan Kuantor}

\begin{explain}
  Teorema kelengkapan juga memerlukan agar deduksi aksiomatik menghasilkan
  fakta-fakta tentang~$\Proves$ yang dibuktikan dalam bagian ini.
\end{explain}

\begin{thm}
\ollabel{thm:strong-generalization} Jika $c$ adalah !!a{constant} yang
tidak muncul dalam $\Gamma$ ataupun $!A(x)$ dan $\Gamma \Proves !A(c)$,
maka $\Gamma \Proves \lforall[x][!A(x)]$.
\end{thm}

\begin{proof}
Menurut teorema deduksi, $\Gamma \Proves \ltrue \lif !A(c)$. Karena
$c$ tidak muncul dalam $\Gamma$, $\top$, ataupun~$!A(x)$, kita memperoleh
$\Gamma \Proves \ltrue \lif \lforall[x][!A(x)]$. Karena aksioma
kebenaran tersedia, modus ponens selanjutnya memberikan
$\Gamma \Proves \lforall[x][!A(x)]$.
\end{proof}

\begin{prop}
\ollabel{prop:provability-quantifiers}
Untuk setiap suku tertutup~$t$:
\begin{tagenumerate}{prvEx,prvAll}
\tagitem{prvEx}{$!A(t) \Proves \lexists[x][!A(x)]$.}{}

\tagitem{prvAll}{$\lforall[x][!A(x)] \Proves !A(t)$.}{}
\end{tagenumerate}
\end{prop}

\begin{proof}
\begin{tagenumerate}{prvEx,prvAll}
\tagitem{prvEx}{Menurut \olref[qua]{ax:q2} dan teorema deduksi.}{}
\tagitem{prvAll}{Menurut \olref[qua]{ax:q1} dan teorema deduksi.}{}
\end{tagenumerate}
\end{proof}

\end{document}
